\documentclass[11pt,a4paper]{article}

\usepackage{amsmath,amssymb}
\usepackage{mathtools}
\usepackage{hyperref}
\usepackage{booktabs}
\usepackage[margin=1.15in]{geometry}

% Shortcuts
\newcommand{\Gstar}{G^{*}}
\newcommand{\vpi}{\varpi}
\DeclareMathOperator{\csch}{csch}

\title{\textsc{One Unit of Existence}\\[6pt]
\large The Life and Death of a Hydrogen Atom\\
as Told by the Constants That Made It}

\author{}
\date{}

\begin{document}
\maketitle
\thispagestyle{empty}

\begin{quote}\itshape
``What is a hydrogen atom?'' is a question that admits an answer at every
level of description.  This is the answer at the level where the description
itself becomes the thing described.
\end{quote}

%=======================================================================
\section*{I.\quad Before the Beginning}
%=======================================================================

Before there is a hydrogen atom, there is a lattice.  Not a lattice made of
anything, but the lattice that makes anything possible.

The lemniscate constant $\vpi\approx 2.622$ is the half-perimeter of the
figure-eight curve $r^2=\cos 2\theta$.  It measures the intrinsic geometry of
the simplest closed curve that crosses itself: the minimal topology of
self-relation.  From~$\vpi$, a second constant is born---the bridge constant
$\Gstar = 2\vpi/\sqrt{\pi}\approx 2.959$---and from these two alone, the
ratio
\[
\lambda \;=\; \frac{\vpi^2}{\Gstar^2} \;=\; \frac{\pi}{4}
\]
generates everything that follows.  In the FTD framework, the number~$\pi$
is not treated as fundamental.  It is $4\lambda$, four times the ratio of
self-intersection to bridging.  This is the \emph{ontic ratio}: the single
number from which the architecture of discreteness and continuity will be
built.

Nothing has happened yet.  No particles, no fields, no space.  Only the
ratio~$\lambda$ and the integers $\{3,4,7,13\}$ that will select, from the
continuum of mathematical possibility, one specific universe.

%=======================================================================
\section*{II.\quad The Quark Condensate: Genesis of the Proton}
%=======================================================================

The universe cools.  At roughly $10^{-5}$~seconds after whatever initiated
the expansion, the temperature drops below $\sim 150$~MeV---the QCD
confinement scale.  Above this energy, quarks move freely in a plasma.
Below it, they do not.  The transition is a phase change in the strong
force, and it is absolute: there will never again, in the entire subsequent
history of the cosmos, be a free quark at thermal equilibrium.

Three quarks---two up, one down---are confined into a proton.  The
confinement is not a cage imposed from outside.  It is a topological fact
about the SU(3) color field: the flux tubes between quarks carry energy
proportional to their length, and stretching them apart costs more energy
than creating new quark-antiquark pairs from the vacuum.  The system chooses
the lowest-energy configuration, which is a bound state.

The number of colors is not arbitrary.  In FTD, the master quadratic
\[
x^2 - 16\,\Gstar^2\,x + 16\,\Gstar^3 = 0
\]
has two roots.  The larger root gives $1/\alpha$.  The smaller root,
$x_-\approx 3.024$, has integer part~3: this is $N_c$, the number of color
charges.  The one-loop QCD beta-function coefficient
$b_3=(11N_c-2N_f)/3=7$ (for $N_f=6$ quark flavors) then sets the running
of the strong coupling.  The proton's confinement is, in this reading,
controlled by the same quadratic that will later determine its
electromagnetic binding to an electron.

The proton's mass is approximately 938~MeV\@.  Almost none of this comes
from the quarks themselves (the up and down quarks have masses of a few~MeV
each).  The overwhelming majority---over 99\%---is the kinetic and binding
energy of the strong field, whose confinement scale
$\Lambda_{\text{QCD}}$ is set by $\alpha_s$ running under $b_3=7$.
The proton is, in this precise sense, made of motion, not of matter.

This is relevant.  The proton is already a discrete-to-continuous bridge: a
discrete, countable object (baryon number $=1$) whose internal structure is
a continuous quantum field theory.  The proton will outlive every star.  Its
lifetime exceeds $10^{34}$~years~\cite{superK}.  For the purposes of this
story, it is eternal.

%=======================================================================
\section*{III.\quad Recombination: The Electron Finds Its Home}
%=======================================================================

For the first 380,000~years, the universe is a plasma.  Protons exist,
electrons exist, photons exist, but the temperature is too high for neutral
atoms to form.  Every time an electron settles into orbit around a proton, a
photon with energy above 13.6~eV strikes it free.  The universe is opaque:
photons scatter constantly off the dense sea of charged particles.

Then the temperature drops below $\sim 3000$~K, and the photon bath no
longer carries enough energy to ionize hydrogen.  In a cosmic eyeblink,
nearly every proton captures an electron.  The universe becomes transparent.

This is the moment our hydrogen atom is born.

The electron does not orbit the proton like a planet.  It occupies a
standing wave, a stationary state of the Schr\"odinger equation:
\[
E_n = -\,\frac{\alpha^2\,m_e c^2}{2n^2}\,,
\qquad n=1,2,3,\ldots
\]
where $\alpha$ is the fine-structure constant.

Here is where the story turns.  The fine-structure constant is the number
that determines the strength of the electromagnetic interaction---the force
that holds our atom together.  In the FTD framework, it emerges from the
master quadratic whose coefficients are built entirely from~$\Gstar$ (and
therefore from~$\vpi$), yielding
\[
\frac{1}{\alpha}\approx 137.036\,,
\]
at a precision of 0.21~parts per trillion~\cite{codata}.  The four integers
$\{3,4,7,13\}$ select this quadratic from the space of all possible
quadratics over the lemniscatic lattice.  The hydrogen atom's entire energy
spectrum---every line in the Balmer series, every transition that Bohr,
Sommerfeld, and Dirac painstakingly catalogued---follows from this single
algebraic object.

As a concrete checkpoint: the ground-state binding energy is
\[
|E_1| = \frac{\alpha^2\,m_e c^2}{2}
= \frac{(1/137.036)^2\times 0.51100\;\text{MeV}}{2}
= 13.606\;\text{eV}\,,
\]
matching the measured ionization energy of hydrogen to five
significant figures.  This is not a fit---$\alpha$ and $m_e$ are both
determined within FTD; the 13.6~eV follows as output.

The electron settles into the $n=1$ ground state.  Its ``orbital'' is a
spherically symmetric cloud, densest at the nucleus, decaying exponentially
with a characteristic scale set by the Bohr radius:
\[
a_0 = \frac{\hbar}{\alpha\,m_e c}
= \frac{1}{\alpha}\cdot\frac{\hbar}{m_e c}
\approx 0.529\;\text{\AA}\,.
\]
Notice: $a_0$ is the Compton wavelength of the electron divided by~$\alpha$.
The Bohr radius is not an independent parameter.  It is $\alpha$'s signature
in configuration space, and $\alpha$ is $\Gstar$'s signature in coupling
space.

The atom now exists.  It is one proton, one electron, and a quantum field
that binds them.  It weighs slightly less than its parts: the mass deficit
is the binding energy, $13.6\;\text{eV}/c^2$, a tiny fraction of the total
938~MeV\@.  It is electrically neutral.  It is the simplest bound state in
nature.

%=======================================================================
\section*{IV.\quad The Discrete Spectrum: A Sum over Integers}
%=======================================================================

Our hydrogen atom sits in its ground state.  It is silent, invisible, inert.
Then a photon arrives.

If the photon carries exactly the right energy---$E_2-E_1=10.2$~eV---the
atom absorbs it and the electron jumps to $n=2$.  Not approximately.
Exactly.  The atom either absorbs the photon or it does not.  There is no
half-transition.  Discreteness is absolute.

The photon energies that the atom can absorb form a discrete set:
\[
E_{n\to m} = \frac{\alpha^2\,m_e c^2}{2}
\left(\frac{1}{n^2}-\frac{1}{m^2}\right),
\qquad n<m.
\]
This is a sum over integers.  The same \emph{structural pattern}---discrete
indices bridged to a continuous closed form by an analytic
continuation---appears in the identity
\[
\sum_{n=0}^{\infty}\frac{1}{(n^2+1)^2}
= \frac{2+\pi\coth\pi+\pi^2\csch^2\pi}{4}
= \frac{2+4\lambda\coth(4\lambda)+16\lambda^2\csch^2(4\lambda)}{4}\,.
\]
The analogy is one of \emph{form}, not of mechanism: the Rydberg series
arises from the Coulomb Green's function, while the sum above passes
through the Gaussian lattice $\mathbb{Z}[i]$.  But the architectural
motif is the same.  Discrete energy levels are the integers of the atom's
inner life.  The electromagnetic field in which they are embedded is the
continuous medium.  The coupling constant~$\alpha$---determined
by~$\Gstar$---is the bridge between them.

The hydrogen emission spectrum, recorded on a photographic plate, is a set
of sharp lines: discrete, countable, unmistakable.  But the underlying field
theory that produces those lines (quantum electrodynamics) is continuous,
built on integrals over all of momentum space.  The sharp lines emerge from
the continuum exactly as the closed-form answer emerges from the infinite
discrete sum---different mechanisms, same structural crossing from
$\sum$ to~$\int$.

In both cases, the bridge is the same constant in different clothing.

%=======================================================================
\section*{V.\quad The Lamb Shift: Where the Vacuum Speaks}
%=======================================================================

If the hydrogen atom were governed by the Dirac equation alone, the
$2S_{1/2}$ and $2P_{1/2}$ states would have exactly the same energy.  They
do not.  The $2S_{1/2}$ state is shifted upward by about 1058~MHz---the Lamb
shift, measured by Willis Lamb in 1947~\cite{lamb47}.

The shift comes from the vacuum.  Even in its ground state, the
electromagnetic field fluctuates: virtual photons are constantly created and
annihilated, and the electron interacts with them.  These vacuum
fluctuations slightly smear the electron's position, effectively making it
``see'' a slightly different nuclear potential than it would classically.
The $S$ states, which have nonzero probability density at the nucleus, are
more affected than $P$ states, which vanish there.

The leading-order Lamb shift scales as $\alpha^5$: five powers of the
fine-structure constant, five powers of $\Gstar$'s imprint.  Each power
of~$\alpha$ represents one more layer of interaction between the discrete
(the atom's energy levels) and the continuous (the quantum vacuum).  The
Lamb shift is the discrete-to-continuous bridge made measurable: the atom's
sharp energy levels responding to the hum of the infinite continuum beneath
them.

This was the measurement that confirmed quantum electrodynamics.  The vacuum
is not empty.  The vacuum is the continuum from which discreteness is
projected.

%=======================================================================
\section*{VI.\quad The Hyperfine Transition: The Quietest Sound}
%=======================================================================

The proton has spin.  The electron has spin.  When aligned, the total energy
is slightly higher than when anti-aligned.  The difference is the hyperfine
splitting: $5.9\;\mu$eV, corresponding to a photon frequency of
1420.405~MHz, or a wavelength of 21~cm.

This is the most famous frequency in radio astronomy.  Neutral hydrogen
throughout the universe occasionally undergoes the spin-flip transition,
emitting a 21-cm photon.  The transition is fantastically slow: the
spontaneous lifetime is about $10^7$~years.  A single hydrogen atom in its
excited hyperfine state will, on average, wait ten million years before
emitting.

But there are $\sim 10^{80}$ hydrogen atoms in the observable universe.  So
the 21-cm line is everywhere, all the time, in every direction.  It is the
sound of the universe breathing.

The hyperfine splitting scales as $\alpha^4\,(m_e/m_p)$: the fine-structure
constant to the fourth power, modulated by the electron-to-proton mass
ratio.  The mass ratio introduces the strong force (since $m_p$ is set by
QCD, whose confinement scale runs under $b_3=7$), so the hyperfine line is a
meeting point of two of the fundamental interactions, both encoded, in the
FTD picture, in the algebraic structure of the lemniscatic lattice.

%=======================================================================
\section*{VII.\quad Stellar Fusion: Death by Continuity}
%=======================================================================

Our hydrogen atom drifts for hundreds of millions of years, part of a
diffuse cloud in intergalactic space.  Then gravity gathers the cloud.  It
contracts, heats, and eventually the core temperature reaches $\sim 10^7$~K.

At this temperature, protons have enough kinetic energy to overcome their
mutual Coulomb repulsion and approach within $\sim 1$~fm of each other,
where the strong nuclear force can act.  Two protons fuse: one undergoes
beta-plus decay (converting to a neutron via the weak force), and the result
is a deuteron, a positron, and a neutrino.

Our hydrogen atom's proton has just become part of a deuterium nucleus.  The
atom, as a distinct entity, is gone.

The probability of this fusion event is controlled by quantum tunneling
through the Coulomb barrier.  The tunneling rate depends exponentially on the
Gamow energy:
\[
E_G = 2\,m_r c^2\,(\pi\alpha Z_1 Z_2)^2\,,
\]
where $\alpha$ appears again.  The exponential sensitivity means that even a
small change in~$\alpha$ would either extinguish all stars or cause them to
burn through their fuel in moments.  The value of~$\alpha$ implied by the
FTD master quadratic happens to fall in the narrow window compatible with
stellar lifetimes of $\sim 10^{10}$~years, long enough for complexity to
emerge on their planets.

This is not a teleological claim.  It is an observational constraint.
Whether the compatibility of FTD's~$\alpha$ with stellar longevity
reflects selection, coincidence, or necessity is an open question.  FTD
provides the value; it does not (yet) explain why the value is
life-permitting.

%=======================================================================
\section*{VIII.\quad The Fusion Chain: From Hydrogen to Iron}
%=======================================================================

The deuteron that was once our atom's proton captures another proton to form
helium-3, and two helium-3 nuclei fuse to form helium-4 plus two protons
returned to the plasma.  This is the pp~chain, the dominant energy source in
solar-type stars.

The mass of helium-4 is less than the mass of four protons by 26.7~MeV\@.
This mass deficit, converted to energy via $E=mc^2$, is what powers the
star.  The sun converts about 600~million tonnes of hydrogen into helium
every second.  Our atom's proton contributed its share: 6.7~MeV per nucleon,
released as photons and neutrinos.

In more massive stars, the chain continues: helium to carbon (via the
triple-alpha process), carbon to oxygen, oxygen to neon, neon to silicon,
silicon to iron.  At iron, fusion becomes endothermic.  The chain stops.
The star has converted the maximum possible fraction of its rest mass into
binding energy, and the nucleus of iron-56 sits at the bottom of the valley
of stability.

Each step in this chain is a discrete nuclear transition, but the stellar
interior in which it occurs is a continuous plasma.  The rates of these
transitions are set by coupling constants (strong, weak, electromagnetic)
that are, in the FTD picture, determined by the integers $\{3,4,7,13\}$
acting on the lemniscatic lattice.  The entire periodic table is a
consequence of these four numbers.

%=======================================================================
\section*{IX.\quad Supernova: The Return}
%=======================================================================

If the star is massive enough, the iron core collapses.  In less than a
second, a volume the size of the Earth is compressed to a sphere 20~km
across.  The collapse is halted by neutron degeneracy pressure---a quantum
effect, the Pauli exclusion principle forbidding two neutrons from occupying
the same state.  The infalling matter bounces off this incompressible core
and an outward shock wave tears the star apart.

This is a supernova.  It releases $\sim 3\times 10^{46}$~J in a few
seconds, briefly outshining its entire host galaxy.  The shock wave
disperses the star's outer layers---including all the elements heavier than
helium that were forged in the core---back into the interstellar medium.

Among the debris: new hydrogen atoms, formed as the expanding ejecta cool
and free protons capture electrons from the thinning plasma.  The cycle is
ready to begin again.

But our particular proton is likely buried deep, fused into iron or heavier
elements synthesized in the explosion's neutron flux.  It is no longer
hydrogen.  Its identity as a distinct baryon persists (baryon number is
conserved, as far as any experiment has ever shown), but its context has
changed irreversibly.

%=======================================================================
\section*{X.\quad Annihilation: The Final Bridge}
%=======================================================================

There is one more chapter, though it may never be written.

If baryon number is not exactly conserved---if grand unified theories are
correct and the proton eventually decays---then the last trace of our
hydrogen atom will vanish.  Many GUT models predict the dominant channel
$p\to e^+ + \pi^0$, followed by $\pi^0\to 2\gamma$ and $e^+$ annihilating
with an ambient electron to produce two more photons.  The entire rest mass
of the proton is converted to radiation.  Whether this channel is realized,
and whether FTD's algebraic structure permits or forbids it, remains an open
question.

The proton's lifetime is at least $10^{34}$~years~\cite{superK}, far longer
than the current age of the universe ($1.4\times 10^{10}$~years).  If it
decays at all, it will do so in an era when every star has burned out, every
black hole has evaporated, and the universe consists of nothing but a thin,
cooling gas of photons, neutrinos, and the occasional lonely electron or
positron.

At that moment, the last discrete object---the proton, baryon number
one---crosses the final bridge into the continuum of radiation.  Photons
have no rest mass, no charge, no baryon number.  They are pure field, pure
wave, pure continuity.  The integer becomes the irrational.  The lattice
point dissolves into the plane.

\[
\sum \;\longrightarrow\; \int
\]

%=======================================================================
\section*{Coda:\quad The Ratio Remains}
%=======================================================================

Every stage of this story was governed by the same architecture:

\bigskip
\begin{center}
\begin{tabular}{lll}
\toprule
\textbf{Event} & \textbf{Discrete} & \textbf{Continuous}\\
\midrule
Quark confinement & baryon number & QCD flux tubes\\
Recombination & electron capture & photon field\\
Absorption/emission & energy levels $E_n$ & electromagnetic field\\
Lamb shift & atomic state & vacuum fluctuations\\
Hyperfine transition & spin states & radio-frequency field\\
Stellar fusion & nuclear binding & thermal plasma\\
Proton decay & baryon number $\to 0$ & radiation field\\
\bottomrule
\end{tabular}
\end{center}
\bigskip

At every transition, the bridge between the discrete side and the continuous
side is controlled by coupling constants that are, in the FTD framework,
determined by $\vpi$, $\Gstar$, and the integers $\{3,4,7,13\}$.

The hydrogen atom does not ``know'' about the lemniscatic lattice.  It does
not need to.  It simply \emph{is} what the lattice describes, projected into
three spatial dimensions and one temporal one, bound by a coupling constant
that is a root of a quadratic over~$\Gstar$, radiating at frequencies that
are integer multiples of $\alpha^2 m_e c^2$, and destined to dissolve back
into the continuum from which it was projected.

The ratio $\lambda=\vpi^2/\Gstar^2$ was there before the atom, is there
during the atom, and will be there after the atom.  The atom is a transient
excitation.  The constants are not made of atoms.

The bridge is not $\pi$.  The bridge is $\Gstar$.

And the story of one hydrogen atom is the story of every hydrogen atom,
which is the story of every discrete thing that ever existed in a continuous
universe.

\bigskip
\begin{center}\itshape
Foundational Ternary Dynamics: a framework, not a faith.
\end{center}

%=======================================================================
\begin{thebibliography}{9}

\bibitem{codata}
E.~Tiesinga \textit{et al.},
``CODATA recommended values of the fundamental physical constants: 2022,''
\textit{Rev.\ Mod.\ Phys.}\ \textbf{96} (2024), 015002.

\bibitem{lamb47}
W.\,E.~Lamb and R.\,C.~Retherford,
``Fine structure of the hydrogen atom by a microwave method,''
\textit{Phys.\ Rev.}\ \textbf{72} (1947), 241--243.

\bibitem{superK}
Super-Kamiokande Collaboration,
``Search for proton decay via $p\to e^+\pi^0$ and $p\to\mu^+\pi^0$,''
\textit{Phys.\ Rev.\ D}\ \textbf{95} (2017), 012004.

\end{thebibliography}

\end{document}
